\documentclass[11pt]{article}
\usepackage[margin=0.9in]{geometry}
\usepackage{amsmath,amssymb,amsthm,mathtools}
\usepackage{graphicx}
\usepackage[colorlinks=true,linkcolor=blue,citecolor=blue,urlcolor=blue]{hyperref}
\usepackage{microtype}

\newtheorem{theorem}{Theorem}
\newtheorem{corollary}{Corollary}
\newcommand{\D}{\mathrm D}
\newcommand{\Cov}{\operatorname{Cov}}

\title{Sharp all-eigenvalue Shannon--Stam stability\\
for Gaussian covariance pairs}
\author{Candidate partial-result packet}
\date{11 August 2026}

\begin{document}
\maketitle

\noindent\textbf{Status.}
Candidate substantial partial result; likely valid; human review recommended.
The Poincar\'e-exponent question for arbitrary log-concave inputs remains open
here.

\section{Source question}

Eldan and Mikulincer prove a Shannon--Stam stability estimate for centered
log-concave vectors whose constant deteriorates cubically in a normalized
Poincar\'e parameter~\cite{EM}.  They ask whether the cubic general bound or
the quadratic uniformly log-concave bound is tight.  They also suggest that
the minimum-eigenvalue/Poincar\'e factor should be replaced by a quantity
using all covariance eigenvalues.  The exact source passage (PDF page~4) is
reproduced below.

\begin{center}
\includegraphics[width=0.90\textwidth]{figures/open_problem_crop.png}
\end{center}

The result below completely answers the covariance part on the subclass of
all centered Gaussian pairs.  It shows, in particular, that arbitrarily
ill-conditioned Gaussians cannot force either Poincar\'e exponent.

\section{Proof idea}

Let the average covariance be $C$.  A common whitening sends the two
covariances to $I+A$ and $I-A$.  Since these matrices are polynomials in one
symmetric matrix, every quantity diagonalizes in the eigenbasis of $A$.
The deficit and the total relative entropy become sums of elementary scalar
logarithmic expressions.

For $\lambda\leq 1/2$, put $c=1-2\lambda$.  The lower estimate is exactly the
concavity inequality
\[
 \log(1-ca)\geq c\log(1-a),
 \qquad 1-ca=c(1-a)+(1-c)\cdot1.
\]
The upper estimate follows because
$\log(1-ca)+c\log(1+a)$ has its maximum at $a=0$.  Swapping the inputs proves
the other half of the parameter range.  No spectral estimate or Poincar\'e
inequality is used.

\section{Exact formula and sharp stability}

For a positive definite matrix $C$, write $G_C$ for the centered Gaussian
with covariance $C$.  The Shannon--Stam deficit is
\[
 \delta_{\lambda}(X,Y)
 =\lambda\D(X\|G_C)+(1-\lambda)\D(Y\|G_C)
 -\D\bigl(\sqrt\lambda X+\sqrt{1-\lambda}Y\,\|G_C\bigr).
\]
The right side is independent of the chosen centered Gaussian reference,
because its constant and quadratic terms cancel; it is therefore the same
deficit as in the source.

\begin{theorem}[sharp Gaussian spectral stability]\label{thm:main}
Let $X$ and $Y$ be independent centered Gaussian vectors in $\mathbb R^d$
with positive definite covariance matrices $\Sigma_X$ and $\Sigma_Y$.  Set
\[
 C=\frac{\Sigma_X+\Sigma_Y}{2},\qquad
 A=\frac12 C^{-1/2}(\Sigma_X-\Sigma_Y)C^{-1/2},
\]
and let $a_1,\ldots,a_d\in(-1,1)$ be the eigenvalues of $A$.  For every
$0<\lambda<1$,
\begin{align}
 \delta_\lambda(X,Y)
 &=\frac12\sum_{i=1}^d\bigl[
 \log(1+(2\lambda-1)a_i)-\lambda\log(1+a_i)
 -(1-\lambda)\log(1-a_i)\bigr],\label{eq:deficit}\\
 \D(X\|G_C)+\D(Y\|G_C)
 &=\frac12\sum_{i=1}^d-\log(1-a_i^2).\label{eq:entropy}
\end{align}
Moreover,
\begin{align}
 \min\{\lambda,1-\lambda\}
 [\D(X\|G_C)+\D(Y\|G_C)]
 &\leq \delta_\lambda(X,Y)\label{eq:lower}\\
 &\leq \max\{\lambda,1-\lambda\}
 [\D(X\|G_C)+\D(Y\|G_C)].\label{eq:upper}
\end{align}
Both constants are optimal, already in dimension one.  At $\lambda=1/2$
there is the identity
\[
 \delta_{1/2}(X,Y)
 =\frac12[\D(X\|G_C)+\D(Y\|G_C)]
\]
for every such covariance pair.
\end{theorem}

\begin{proof}
Apply the common invertible map $C^{-1/2}$.  This leaves the entropy deficit
unchanged and sends the covariances to $I+A$ and $I-A$.  The covariance of
$\sqrt\lambda X+\sqrt{1-\lambda}Y$ becomes
$I+(2\lambda-1)A$.

For centered Gaussians,
\[
 \D(G_\Sigma\|G_I)
 =\frac12\{\operatorname{tr}\Sigma-d-\log\det\Sigma\}.
\]
The trace terms cancel in the deficit.  Diagonalizing $A$ gives
\eqref{eq:deficit}.  Adding the two input entropies also cancels the trace
terms and gives
\[
 -\frac12\log\det((I+A)(I-A))
 =\frac12\sum_i-\log(1-a_i^2),
\]
which is \eqref{eq:entropy}.

It remains to prove a scalar statement.  Define
\begin{align*}
 f_\lambda(a)&=\log(1+(2\lambda-1)a)
 -\lambda\log(1+a)-(1-\lambda)\log(1-a),\\
 d(a)&=-\log(1-a^2).
\end{align*}
Suppose first that $0<\lambda\leq1/2$ and put $c=1-2\lambda\in[0,1)$.
Then
\[
 f_\lambda(a)-\lambda d(a)
 =\log(1-ca)-c\log(1-a)\geq0
\]
by concavity of the logarithm and
$1-ca=c(1-a)+(1-c)\cdot1$.  On the other hand,
\[
 f_\lambda(a)-(1-\lambda)d(a)
 =h_c(a):=\log(1-ca)+c\log(1+a).
\]
Here
\[
 h_c'(a)=\frac{-c(1+c)a}{(1-ca)(1+a)},
\]
so $h_c$ attains its maximum $h_c(0)=0$ at zero.  Hence
$\lambda d(a)\leq f_\lambda(a)\leq(1-\lambda)d(a)$.
For $\lambda\geq1/2$, use
$f_\lambda(a)=f_{1-\lambda}(-a)$ and $d(-a)=d(a)$.  Summing the scalar
inequalities proves \eqref{eq:lower}--\eqref{eq:upper}.

If $\lambda<1/2$, the ratio $f_\lambda(a)/d(a)$ tends to $\lambda$ as
$a\downarrow-1$ and to $1-\lambda$ as $a\uparrow1$.  The limits are reversed
when $\lambda>1/2$.  Thus both constants are sharp in dimension one.  When
$\lambda=1/2$, the scalar lower and upper bounds coincide, proving the stated
identity.
\end{proof}

\begin{corollary}[source normalization]
If $\Sigma_X+\Sigma_Y=2I$, then $C=I$ and Theorem~\ref{thm:main} uses exactly
the source's standard Gaussian $G$.  In particular,
\[
 \delta_{EPI,\lambda}(X,Y)\geq
 \min\{\lambda,1-\lambda\}\,[\D(X\|G)+\D(Y\|G)],
\]
uniformly in the dimension, the smallest covariance eigenvalue, and the
Gaussian Poincar\'e constants.
\end{corollary}

\section{Scope and obstruction to a full upgrade}

The theorem treats every positive definite Gaussian covariance pair, not
only commuting covariances assumed in advance: after average whitening they
automatically have the form $I\pm A$.  Equations~\eqref{eq:deficit} and
\eqref{eq:entropy} use every eigenvalue and remain sharp as one eigenvalue
tends to degeneracy.

For non-Gaussian inputs, entropy bookkeeping separates the deficit into the
Gaussian covariance deficit plus
\[
 \lambda\D(X\|G_X)+(1-\lambda)\D(Y\|G_Y)-\D(Z\|G_Z),
\]
where each $G_\bullet$ matches the covariance of its variable and
$Z=\sqrt\lambda X+\sqrt{1-\lambda}Y$.  This residual has no suitable sign
from ordinary convolution data processing when the input covariances differ.
Eight focused routes---including perturbation, heat flow, same-law
polarization, anisotropic preconditioning, and product extremizers---all
return to control of this nonlinear two-law interaction.  They are documented
in the associated attempt note.

A bounded novelty search on 11 August 2026 used the exact arXiv identifier and
question, Gaussian Shannon--Stam covariance stability, log-determinant
formulas, Poincar\'e-exponent phrases, and later equality/stability work.  It
found the source and adjacent qualitative or regularized results, but no exact
match for Theorem~\ref{thm:main} and no later resolution of the general
exponent question.  Because the Gaussian argument is elementary and may be
implicit in matrix-entropy literature, novelty confidence is moderate.

\paragraph{Human-review recommendation.}
Audit the reference-Gaussian cancellation and the two scalar logarithmic
inequalities.  These are the only substantive steps.  Also check whether the
sharp two-sided Gaussian formula is already recorded under a different
log-determinant Jensen-gap terminology.

\begin{thebibliography}{9}
\bibitem{EM}
R.~Eldan and D.~Mikulincer,
\emph{Stability of the Shannon--Stam inequality via the F\"ollmer process},
ALEA Lat. Am. J. Probab. Math. Stat. \textbf{17} (2020), 487--507;
arXiv:1903.07140.
\end{thebibliography}

\end{document}
